% Image credits for the photographs and AI illustrations of Book 3
% (University Chemistry, Year 2). Machine-readable license records live in
% images/book3/CREDITS.md; keep the two files in sync. The Book 3 writer
% owns this file: add one \item per photograph (in an itemize, after the
% first paragraph) as photographs land.
\chapter*{Kredit Gambar}

Gambar-gambar dalam buku ini adalah karya asli. Foto-foto, jika dipakai,
dipakai dengan lisensi bebasnya masing-masing, dengan terima kasih kepada
para pembuatnya; tautan sumber dan URL lisensi lengkap untuk setiap foto
tercantum dalam \texttt{images/book3/CREDITS.md} di repositori buku ini.
Piktogram bahaya adalah piktogram Sistem Harmonisasi Global untuk
Klasifikasi dan Pelabelan Bahan Kimia, sebagaimana digambar untuk paket
\texttt{ghsystem} dari \TeX\ Live.

\bigskip
Di samping foto-foto dan gambar-gambar asli, beberapa gambar adalah
ilustrasi buatan AI, yang dihasilkan dengan pembangkit gambar OpenAI dari
perintah yang ditulis oleh para penyunting buku ini, dan diperiksa
ketepatan kimianya sebelum dimuat. Perintah lengkap untuk setiap gambar
yang dihasilkan dicatat dalam \texttt{images/book3/ai/PROMPTS.md} di
repositori.

\bigskip
\noindent Foto-foto, menurut bab:
\begin{itemize}\small
  \item Bab~1: Germain Henri Hess, potret oleh I.~I. Reimers (1837--1838),
        CC BY-SA 4.0; Marcellin Berthelot, pembuat tidak diketahui, domain
        publik (keduanya dari Wikimedia Commons).
  \item Bab~2: Josiah Willard Gibbs, fotografer tidak diketahui, domain
        publik (Wikimedia Commons).
  \item Bab~3: François-Marie Raoult, fotografer tidak diketahui, domain
        publik (Wikimedia Commons).
  \item Bab~4: Jacobus Henricus van 't Hoff, oleh Nicola Perscheid (1904),
        domain publik; Henry Le Chatelier, pembuat tidak diketahui, CC BY-SA
        4.0; Fritz Haber dan Carl Bosch, Yayasan Nobel, domain publik
        (semuanya dari Wikimedia Commons).
  \item Bab~6: tanur tinggi di Duisburg-Nord, oleh Dietmar Rabich, CC BY-SA
        4.0; pengelasan termit sebuah rel, oleh PetrS., CC BY-SA 3.0
        (keduanya dari Wikimedia Commons).
  \item Bab~7: kolom distilasi sebuah kilang, oleh CEphoto, Uwe Aranas,
        CC BY-SA 3.0 (Wikimedia Commons).
  \item Bab~9: Michael Faraday, oleh Thomas Phillips (1842), domain publik
        (Wikimedia Commons).
  \item Bab~10: Jaroslav Heyrovsk\'y, arsip Institut Kimia Fisika
        J.~Heyrovsk\'y, CC BY-SA 3.0 (Wikimedia Commons).
  \item Bab~11: sebuah tumpukan Volta, oleh GuidoB, CC BY-SA 3.0; Charles
        Martin Hall dan Paul H\'eroult, pembuat tidak diketahui, domain
        publik (semuanya dari Wikimedia Commons).
  \item Bab~12: sebuah patung tembaga, oleh Elcobbola, domain publik;
        sebuah anode korban, oleh Zwergelstern, CC BY-SA 3.0 (keduanya dari
        Wikimedia Commons).
  \item Bab~13: Erwin Schr\"odinger, Yayasan Nobel (1933), domain publik
        (Wikimedia Commons).
  \item Bab~14: oksigen cair di dalam magnet, oleh Pieter Kuiper, domain
        publik; Robert Mulliken dan Friedrich Hund, oleh GFHund, CC BY 3.0
        (keduanya dari Wikimedia Commons).
  \item Bab~16: August Kekul\'e (1873), pembuat tidak diketahui, domain
        publik (Wikimedia Commons).
  \item Bab~18: larutan-larutan logam transisi, oleh Benjah-bmm27, domain
        publik; Alfred Werner, Les Prix Nobel (1914), domain publik
        (keduanya dari Wikimedia Commons).
  \item Bab~19: sebuah kristal rubi, domain publik; sebuah zamrud, oleh
        Didier Descouens, CC BY-SA 3.0; kristal tembaga sulfat, oleh
        Stephanb, CC BY-SA 3.0 (semuanya dari Wikimedia Commons).
  \item Bab~20: Akira Suzuki, Ei-ichi Negishi, dan Richard Heck, oleh
        BloodIce, CC BY-SA 4.0 (Wikimedia Commons).
  \item Bab~22: Charles Friedel (1890-an) dan James Mason Crafts
        (\textit{Popular Science Monthly}, 1900), domain publik (Wikimedia
        Commons).
  \item Bab~24: Michel Eug\`ene Chevreul, oleh Nicolas Eustache Maurin,
        domain publik (Wikimedia Commons).
  \item Bab~25: Alexander Borodin, pembuat tidak diketahui, domain publik;
        Ludwig Claisen, oleh Veronica.villagrasa, CC BY-SA 4.0 (keduanya
        dari Wikimedia Commons).
  \item Bab~26: Robert Robinson, pembuat tidak diketahui, domain publik
        (Wikimedia Commons).
  \item Bab~28: Elias James Corey, oleh Trvthchem, CC0 (Wikimedia Commons).
  \item Bab~29: Wallace Carothers di laboratorium, fotografer tidak
        diketahui, domain publik (Wikimedia Commons).
  \item Bab~30: Emil Fischer, Atelier Victoria, Berlin, sekitar 1895,
        domain publik (Wikimedia Commons).
  \item Bab~31: Mikhail Tsvet, pembuat tidak diketahui, domain publik
        (Wikimedia Commons).
  \item Bab~32: Francis William Aston (1922), pembuat tidak diketahui,
        domain publik; Gustav Kirchhoff, Robert Bunsen, dan Henry Roscoe
        (1862), domain publik (keduanya dari Wikimedia Commons).
  \item Bab~34: William Sealy Gosset (1908), domain publik (Wikimedia
        Commons).
%% next chapter here
\end{itemize}
\clearpage
